\section{Conclusion}
\label{sec:conclusion}

MAM-AI shows that a fully offline, source-cited medical question-answering
assistant can run on a commodity Android device for nurse-midwives in a
low-connectivity setting. Its layered evaluation relocates the hard problem.
On-device retrieval is essentially solved: on the mamaretrieval benchmark a 300M
embedder ranks third of seven and rivals cloud systems, so the passages the
system needs are usually found. The small generator is what remains in
doubt---at 4B it cannot be both helpful and safe at once: Gemma~3n answers more
fully but commits genuine dangerous errors, while Gemma~4 is safe but less
helpful, whereas the 397B frontier model escapes the trade-off entirely, marking
the gap as a limit of small-model capacity rather than of the task. Corpus
quality is decisive for the same reason: a small model carries little parametric
medical knowledge and leans on what it retrieves, so it needs official,
traceable guidelines and locally relevant protocols---and where the corpus holds
the right passage the answer is specific and safe to act on, where it does not
the answer goes vague (Figure~\ref{fig:case-nausea}).

Generator capability matters---it sets the quality ceiling---but in a
safety-critical setting we prioritise faithfulness, deploying the model that
grounds its answers most reliably and recovering helpfulness through the prompt.
MAM-AI is a thoroughly evaluated research prototype, not a fielded product;
closing the gap to deployment means grounding the generator, improving the
corpus, extending to Swahili, and, above all, testing with the nurse-midwives it
is built for.
